% Nouveau EGA V (cote 25): archive cote 25, site page 38 (PDF page 39).
% typescript. Draft compared: none.
% AI-transcribed by Claude Opus 5.5 (claude-opus-5-5, adaptive thinking, effort high), Message Batches API, batch ega25b (msgbatch_01A7n5YuXubUPNK1zsepi8RG), 2026-11-13; not yet checked by eye.
\scanpage{25/38}
\mspage{\hand{34}}

\margin{$|$}L'ensemble sous-jacent à ce préschéma n'est autre que l'ensemble des
points de $\underline{\hand{\check{}}Y}$= V($\phi$) qui sont \uline{zéros singuliers}\inserted{(par.16)} de la section $\phi$ de
$\underline{O}_{X\mathrm{x}\hand{\check{}}P}$(1,1) \uline{relativement à la base} $\hand{\check{}}$P , i.e. ses points \xout{gé} à valeurs dans
un corps k sur $\hand{\check{}}$P sont les points \inserted{x} de $Y_k\subset X_k$ tels que \xout{\unsure{xx bien} $\phi_k$ \unsure{soit}}
\xout{\unsure{identiquement nulle en x, mais bien}} $\phi_k$ s'annule à l'ordre deux au moins
en x, i.e. tel que \unsure{$\underline{Y}_k$|$\underline{V}_k$} ne soit pas lisse de dimension relative r-1
sur k en x. La caractérisation énoncée des \xout{points} zéros singuliers comm
les éléments du sousèpréschéma lisse P( $\hand{\unsure{\underline{\Omega}}}_{X/P}$) de $X_{\hand{\check{}}P}$ provient aussitôt
de l'énoncé suivant, qui mérite de figurer en proposition préliminaire:
si S=Spec(k), et si H est un hyperplan de P , alors Y = X$\mathrm{x}_P$H est
\xout{x \unsure{sur k} xxxxxxxxxxx}
\xout{lisse en x \unsure{le point rationnel} sur k x si et seulement si} lisse sur k de
dimension relative d-1 en le point x $\in$Y(k) (i.e. x est un zéro non
singulier, i.e. géom non singulier, de la section $\phi$ de $\underline{O}_X$(1) définie
par H), si et seulement si H ne contient pas l'image par $\phi'$ de l'espa-%
ce tangent à X en x (rel à k) - ou \deleted{encore} comme on dit\xout{\illegible{1}} encore
(lorsque f:X$\longrightarrow$P est une immersion, permettant d'identifier X à un
sous-préschéma de P) sss H \uline{n'est pas tangent à} X \uline{en} x. Cela résulte
trivialement du critère jacobien de lissité, ou de la définition d\unsure{u}
zéro singulier, une fois qu'on a bien précisé le sens de l'énoncé,
c'est à dire qu'on a précisé comment un \xout{espace tangent} sous-espace
\margin{$|$}vectoriel de l'espace tangent à P en un point a (=f(x)) définit une
sous-variété linéaire de P (de telle façon que cela \inserted{a} un sens de dire
que H ne contient pas ledit sous-espace vectoriel): cela provient
bien entendu de la suite exacte (b) ci-dessus, qui permet de\xout{identifier}
définir une correspondance biunivoque entre l'ensemble des sous-espaces
vectoriels de l'espace tangent en a, et les sous-variétés linéaires de
P contenant a. Cette correspondance d'ailleurs revient à associer à

\margin{38}
